\section{Out-of-Sample Evidence and Model Risk}

\noindent\begin{minipage}{\columnwidth}
{\centering
  \includegraphics[width=0.95\columnwidth]{figures/oos_rmse.png}
  \captionof{figure}{Dense-only rolling evidence on 30 forecast origins and 4,667 common target observations. The original fits give Bates--Hawkes the lowest mean daily IV RMSE.}
  \label{fig:oos_r2}
\par}
\end{minipage}\par

\noindent\begin{minipage}{\columnwidth}
{\centering
  \includegraphics[width=0.95\columnwidth]{figures/oos_r2.png}
  \captionof*{figure}{\textbf{Figure~\ref{fig:oos_r2} (continued).} Positive $R^2_{OOS}$ indicates improvement over the stated benchmark.}
\par}
\end{minipage}\par

\subsection{Conditional next-date repricing}
Model comparisons use 4,667 targets; persistence comparisons use the 4,026 inside the origin interpolation domain. Benchmark scores therefore have different supports; equal date weighting limits the influence of densely populated dates.

Reported RMSE is the mean of daily RMSEs, not the root of mean daily MSE used by $R^2$. Values are 140.6528 bp for Black--Scholes, 80.1964 bp for Heston, 69.3055 bp for Bates and 65.0507 bp for Full Bates--Hawkes. Relative to the origin mean-IV benchmark, the corresponding $R^2_{OOS}$ values are $-2.58\%$, $66.07\%$, $72.97\%$ and $75.81\%$. None beats interpolated IV-surface persistence, but Full Bates--Hawkes is closest at $-3.68\%$, followed by Bates at $-12.92\%$ and Heston at $-27.24\%$.

\begin{findingbox}[Model-risk conclusion]
Historical OOS uses the original fits; audited current fits are a separate experiment. A paired five-date block bootstrap gives a Hawkes-minus-Bates MSE interval $[-8.95,-2.77]\times10^{-6}$, exploratory with only 30 dates.
\end{findingbox}

\noindent\begin{minipage}{\columnwidth}
{\centering
  \includegraphics[width=0.95\columnwidth]{figures/online_welch_goyal_cumulative.png}
  \captionof{figure}{Cumulative rolling MSE differences against the origin mean IV and interpolated IV-surface persistence. Upward movement favours the option model; the stronger persistence benchmark remains difficult to beat.}
  \label{fig:welch_goyal}
\par}
\end{minipage}\par

\subsection{Jump-layer interpretation}
The self-exciting extension is motivated by clustered option-implied jump activity. Relative to the constant Poisson baseline, the Hawkes specification permits intensity bursts followed by endogenous decay. The branching ratio and decay parameter have different meanings; fitted risk-neutral self-excitation is not evidence of physical news clustering.

Without a dated news filtration or parameter-stability evidence, these option-implied bursts cannot be identified as causal responses to geopolitical events.

\noindent\begin{minipage}{\columnwidth}
{\centering
  \includegraphics[width=0.95\columnwidth]{figures/hawkes_intensity_comparison.png}
  \captionof{figure}{Illustrative simulated Poisson and Hawkes paths under the original parameter set. A single realization cannot establish clustering in observed data.}
  \label{fig:intensity_compare}
\par}
\end{minipage}\par

% Expansion last - 2026-09-06
\subsection{Interpretive limitations and validation}
Reproducibility, numerical accuracy and economic validity are distinct. A good option fit cannot validate the GLD-to-gold mapping or the Q-law/WACC combination. The experiment identifies sensitivity to a chosen numerical contract, not corporate fair value.

The old run equated enterprise and equity proxies using zero bridge placeholders. This revision removes that share-price interpretation; debt, cash and other claims remain unreconciled. The new copper operating extension does not close this corporate bridge. A genuine valuation must start from one dated filing and the matching listed security, establish ownership and attributable production, and distinguish component margin from reported operating profit. Depreciation, cash taxes, working capital, sustaining capital expenditure and growth investment must then enter consistently, with no deduction counted twice. Net debt, non-operating assets, streams and other claims must be matched to a diluted share count on the same basis. Only then can a per-share output acquire a filing-reconciled interpretation.

\subsection{Operating uncertainty and research priorities}
Production, grade, recovery, mine mix and costs may depend on gold prices, inflation and exchange rates. Independence is not automatic: operating dependence and Goodman variance propagation require explicit assumptions. Technology scenarios must include implementation costs and investment; a smooth conditional cost mean does not imply absence of uncertainty.

A mining-company terminal block needs an explicit resource interpretation. Continued production may require reserve conversion, replacement investment or new projects, while depletion and closure liabilities can reduce the duration and value of the cash-flow stream. A positive WACC-growth spread controls a mathematical denominator but does not establish that a perpetual operating horizon is economically justified. Finite-resource and going-concern closures should be compared under coherent investment assumptions before attributing differences entirely to the gold-price model. Such a comparison is a future validation exercise, not a result of the present run.

Market maturities end at 653 days, well before five scenario years and perpetuity. Fixed parameters beyond that horizon are an extrapolation assumption, not market-identified information. The five-year-only comparison removes terminal continuation but does not calibrate depletion, replacement investment or closure costs. A bucketed term-structure and finite-reserve model remain necessary for an economic valuation.

\section*{Conclusion}

Holding operations and corporate assumptions fixed isolates the effect of gold-price model choice. Nested-candidate safeguards improve the current Bates and Hawkes fits. Historical OOS remains distinct and no original engine beats interpolated persistence. Numerical and terminal sensitivity prevent a corporate value ranking.

The copper extension admits three producing mines and adds a conditional five-year operating contribution. Historical snapshots support actual temporal testing, but the sparse rolling experiment does not beat IV persistence. Its positive fixed-origin Bates and Hawkes comparisons are exploratory. The combined output remains an attributable component-margin sensitivity, without a reconciled FCFF or equity bridge.

The persistent-cost assumption is prudential only relative to otherwise identical efficiency-improvement scenarios, not a lower bound. The contribution is an auditable interdisciplinary framework, not a certified value: commodity basis, probability measure, operating uncertainty, corporate accounting and finite-resource closure remain the priorities for economic validation.


\begingroup
\setlength{\bibsep}{0.6pt}
\fontsize{6.75}{7.75}\selectfont
% Keep the frozen bibliography in a regular source file.  Overleaf compiles
% the main document with job name `output`, so depending on a generated .bbl
% file with a different name is fragile when the project is imported.
% Frozen bibliography for direct pdfLaTeX compilation on Overleaf.
% The unused bibliography ledger is archived outside the Overleaf project.
\begin{thebibliography}{}

\bibitem[Akaike, 1974]{akaike1974}
Akaike, H. (1974).
\newblock A new look at the statistical model identification.
\newblock {\em IEEE Transactions on Automatic Control}, 19(6):716--723.

\bibitem[Bates, 1996]{bates1996}
Bates, D.~S. (1996).
\newblock Jumps and stochastic volatility: Exchange rate processes implicit in
  deutsche mark options.
\newblock {\em The Review of Financial Studies}, 9(1):69--107.

\bibitem[Damodaran, 2012]{damodaran2012investment}
Damodaran, A. (2012).
\newblock {\em Investment Valuation: Tools and Techniques for Determining the
  Value of Any Asset}.
\newblock Wiley, 3 edition.

\bibitem[Glasserman, 2004]{glasserman2004monte}
Glasserman, P. (2004).
\newblock {\em Monte Carlo Methods in Financial Engineering}.
\newblock Springer.

\bibitem[Hawkes, 1971]{hawkes1971}
Hawkes, A.~G. (1971).
\newblock Spectra of some self-exciting and mutually exciting point processes.
\newblock {\em Biometrika}, 58(1):83--90.

\bibitem[Heston, 1993]{heston1993}
Heston, S.~L. (1993).
\newblock A closed-form solution for options with stochastic volatility with
  applications to bond and currency options.
\newblock {\em The Review of Financial Studies}, 6(2):327--343.

\bibitem[Schwarz, 1978]{schwarz1978}
Schwarz, G. (1978).
\newblock Estimating the dimension of a model.
\newblock {\em The Annals of Statistics}, 6(2):461--464.


\bibitem[Fang and Oosterlee, 2008]{cos2008}
Fang, F. and Oosterlee, C.~W. (2008). A novel pricing method for European options based on Fourier-cosine series expansions. {\em SIAM J. Sci. Comput.}, 31:826--848.
\bibitem[Svensson, 1994]{nss1994}
Svensson, L. (1994). \href{https://www.nber.org/papers/w4871}{Estimating and interpreting forward interest rates}. NBER WP 4871.
\bibitem[Barrick, 2026]{barrick2026}
Barrick (2026). \href{https://www.barrick.com/English/investors/quarterly-reports/default.aspx}{Q1 and Q2 financial statements, MD\&A and mine statistics}.
\bibitem[Cboe, 2026]{cboe2026}
Cboe (2026). \href{https://www.cboe.com/exchange-traded-stock/etp-options-spec}{ETP options contract specifications}.
\bibitem[SPDR, 2025]{spdr2025}
SPDR (2025). \href{https://files.spdrgoldshares.com/spdr-cms/2025-10/GLD-FAQs-January-2025.pdf}{GLD frequently asked questions}.
\bibitem[U.S. Treasury, 2026]{treasury2026}
U.S. Treasury (2026). \href{https://home.treasury.gov/resource-center-data-chart-center/interest-rates}{Daily Treasury par yield curve rates}, 2 September.

\bibitem[Barrick, 2025]{barrickar2025}
Barrick (2025). \href{https://s25.q4cdn.com/322814910/files/doc_financial/annual_reports/2025/Barrick_Annual_Report_2025.pdf}{Annual Report 2025}, MD\&A and joint-arrangements notes.
\bibitem[Barrick, 2026b]{barrickmine2026}
Barrick (2026). \href{https://s25.q4cdn.com/322814910/files/doc_financial/quarterly_results/2026/q2/Barrick_Q2_2026_Mine_Stats.pdf}{Q2 2026 mine statistics}, copper table, p. 8.

\bibitem[Schwartz, 1997]{schwartz1997}
Schwartz, E.~S. (1997). \href{https://doi.org/10.1111/j.1540-6261.1997.tb02721.x}{The stochastic behavior of commodity prices: implications for valuation and hedging}. {\em Journal of Finance}, 52(3):923--973.

\bibitem[Barrick Gold Research Teams, 2026]{sfcgithub}
Barrick Gold Research Teams. (2026).
\newblock {\em Operating-value sensitivity: code, data and audit tables}.
\newblock Starting Finance Club PoliTo, GitHub repository.
\newblock \href{https://github.com/StartingFinanceClubPoliTo/Research/tree/main/Barrick-Gold/Barrick\%20Mining\%20PAPER}{Barrick Mining research repository}, revised 5 October 2026.

\end{thebibliography}

\endgroup
